\section*{Hoofdstuk \ref{ch:b3:bacteriology} --- Bacteriologie: groei, fysiologie en genetica}

\begin{solution}{exo:b3:bacteriology:1}
\omterm{def:b3:bacteriology:envelope}{Grampositief}: een cytoplasmamembraan gewikkeld in een dikke wand van
\omterm{def:b3:bacteriology:envelope}{peptidoglycaan} uit vele lagen, doorregen met teichoïnezuren.
\omterm{def:b3:bacteriology:envelope}{Gramnegatief}: een dunne laag \omterm{def:b3:bacteriology:envelope}{peptidoglycaan} in een periplasma, daarna
een \omterm{def:b3:bacteriology:envelope}{buitenmembraan} waarvan de buitenste laag uit \omterm{def:b3:bacteriology:envelope}{lipopolysacharide}
bestaat, met porines. Het complex van kristalviolet en jodium blijft
gevangen in de dikke, door alcohol uitgedroogde wand van de
\omterm{def:b3:bacteriology:envelope}{grampositieve} (paars) maar wordt uit de \omterm{def:b3:bacteriology:envelope}{gramnegatieve} gespoeld door
haar dunne wand en haar ontwrichte \omterm{def:b3:bacteriology:envelope}{buitenmembraan}, dat daarna de roze
nakleuring opneemt.
\end{solution}

\begin{solution}{exo:b3:bacteriology:2}
$10^{6}$-voudig $= 2^{19.9}$: ongeveer $20$ generaties; $T = \qty{8}{h}/20
= \qty{24}{min}$; $\mu = \ln 2/\qty{0.4}{h} = \qty{1.7}{h^{-1}}$.
\end{solution}

\begin{solution}{exo:b3:bacteriology:3}
\omterm{def:b3:bacteriology:hgt}{Transformatie} --- Griffith (1928) en Avery, MacLeod en McCarty (1944):
door hitte gedode virulente pneumokokken transformeren levende
avirulente, en het middel is DNA. \omterm{def:b3:bacteriology:hgt}{Transductie} --- Zinder en Lederberg
(1952): faag P22 draagt genen van \emph{Salmonella} van de ene stam
naar de andere door een filter dat cellen tegenhoudt. \omterm{def:b3:bacteriology:hgt}{Conjugatie} ---
Lederberg en Tatum (1946): twee auxotrofe stammen van \emph{E.~coli}
geven alleen prototrofe recombinanten wanneer zij worden gemengd,
waarbij celcontact nodig is (de U-buis van Davis).
\end{solution}

\begin{solution}{exo:b3:bacteriology:4}
Penicilline: de transpeptidasen die het \omterm{def:b3:bacteriology:envelope}{peptidoglycaan} dwarsverbinden,
die dierlijke cellen niet hebben. Streptomycine: de kleine ribosomale
subeenheid, waarvan de bacteriële vorm van de eukaryote verschilt.
Ciprofloxacine: het DNA-gyrase, een bacterieel topo-isomerase dat de
mens niet heeft. Rifampicine: de $\beta$-subeenheid van het bacteriële
RNA-polymerase, die het eukaryote enzym niet deelt.
\end{solution}

\begin{solution}{exo:b3:bacteriology:5}
$S^{*} = K_{S}D/(\mu_{\max} - D)$, $X^{*} = Y(S_{0} - S^{*})$. $D = 0.2$:
$S^{*} = \qty{0.0067}{g/L}$, $X^{*} = \qty{0.80}{g/L}$. $D = 0.6$:
$S^{*} = 0.06$, $X^{*} = 0.78$. $D = 0.79$: $S^{*} = 1.58$, $X^{*} = 0.17$.
Uitspoeling bij $D = \mu(S_{0}) = 0.8\times 2/2.02 = \qty{0.79}{h^{-1}}$.
\end{solution}

\begin{solution}{exo:b3:bacteriology:6}
$(N/V)_{\text{thr}} = kA_{\text{thr}}/p$. Lichtorgaan: $0.01\times
6\times 10^{12}/100 = 6\times 10^{8}$ per liter $= 6\times 10^{5}$ per
mL. Zeewater: $10\times 6\times 10^{12}/100 = 6\times 10^{11}$ per liter
$= 6\times 10^{8}$ per mL --- duizend keer hoger, in zee nooit
gehaald.
\end{solution}

\begin{solution}{exo:b3:bacteriology:7}
\qty{25}{mm}: gevoelig, lage MRC. \qty{12}{mm}: verminderd gevoelig,
intermediair --- kan bij gewone doses falen. \qty{0}{mm}: resistent,
groei tot aan het schijfje. Het middel diffundeert vanaf het schijfje
zodat de concentratie op een reproduceerbare manier met de afstand
daalt; de groei stopt waar de concentratie gelijk is aan de MRC, zodat
de straal van de zone en de MRC door een ijkcurve zijn verbonden die
met stammen van bekende MRC is opgesteld.
\end{solution}

\begin{solution}{exo:b3:bacteriology:8}
Het conjugatieve \omterm{def:b3:genetic-engineering:tools}{plasmide} draagt een integron waarvan de cassettereeks
verscheidene resistentiegenen bevat, plus transposons met nog meer; het
\omterm{def:b3:genetic-engineering:tools}{plasmide} steekt bij één paring over en brengt ze alle tot expressie.
Zij komen bijeen omdat de selectie voor een van de middelen het hele
\omterm{def:b3:genetic-engineering:tools}{plasmide} in stand houdt, zodat genen die erop meeliften blijven bestaan
(co-selectie); integronen vangen cassettes; transposons springen op
\omterm{def:b3:genetic-engineering:tools}{plasmiden}; en het \omterm{def:b3:genetic-engineering:tools}{plasmide}, niet het gen, is de eenheid van overdracht.
\end{solution}

\begin{solution}{exo:b3:bacteriology:9}
Eén middel: $10^{-7}\times 10^{9} = 100$ resistente cellen te
verwachten. Twee onafhankelijke middelen: $10^{-14}\times 10^{9} = 10^{-5}$.
Persisters zijn geen mutanten: een klein deel van de sluimerende cellen
overleeft elk middel en groeit weer uit tot een gevoelige populatie
zodra het middel weg is, zodat een tweede middel niet helpt; alleen een
behandeling die lang genoeg duurt om hen bij het ontwaken te vangen, of
een middel dat op sluimerende cellen werkt, doet dat wel.
\end{solution}

\begin{solution}{exo:b3:bacteriology:10}
$X > X^{*}$: het verbruik $\mu X/Y$ overtreft de aanvoer, dus zakt $S$
onder $S^{*}$; dan is $\mu(S) < D$ en $\mathrm{d}X/\mathrm{d}t < 0$ en
zakt $X$ terug; naarmate $X$ daalt stijgt $S$ weer. Beide afwijkingen
worden gecorrigeerd: de evenwichtstoestand is stabiel. Twee soorten:
elk heeft $\mu_{i}(S) = D$ nodig om te blijven bestaan, wat haar eigen
$S_{i}^{*}$ vastlegt; de voedingsstof kan maar één waarde aannemen,
dus blijft hoogstens één soort bestaan. Die met de laagste $S^{*}(D)$
drukt $S$ onder wat de andere nodig heeft, en de andere wordt
uitgespoeld; welke soort dat is kan met $D$ veranderen als hun
monodcurven elkaar kruisen.
\end{solution}

\begin{solution}{exo:b3:bacteriology:11}
Met dichtheid $\rho = N/V$ bestaat de uittoestand $A = p_{0}\rho/k$
zolang die onder $A_{\text{thr}}$ blijft, dus $\rho < kA_{\text{thr}}/p_{0}$;
de aantoestand $A = p_{1}\rho/k$ bestaat zolang die erboven blijft, dus
$\rho
> kA_{\text{thr}}/p_{1}$. Voor $kA_{\text{thr}}/p_{1} < \rho <
kA_{\text{thr}}/p_{0}$ zijn beide met zichzelf in overeenstemming: een
populatie die is aangegaan houdt zichzelf met haar hogere uitstoot aan,
en een die dat niet is blijft uit. Het nut: binding --- zodra een
kolonie heeft besloten een biofilm te bouwen of toxine vrij te laten,
maakt een dip in de dichtheid dat besluit niet ongedaan, en de
schakeling is scherp en gelijktijdig in plaats van geleidelijk.
\end{solution}

\begin{solution}{exo:b3:bacteriology:12}
(a) Hoge doses gedurende korte kuren: het niveau blijft boven de MRC
van de mutanten en doodt wildtype en mutanten van de eerste stap ---
aan te bevelen, begrensd door de giftigheid. (b) Lage doses gedurende
lange kuren: het niveau zit dagenlang in het venster en kweekt
resistentie --- af te raden. (c) Stoppen zodra de klachten verdwijnen:
de overlevenden zijn de minst gevoelige cellen en het dalende niveau
gaat door het venster --- af te raden. (d) Combinatie: een cel moet
tegen beide bestand zijn, en die bevat een miljard cellen niet --- aan
te bevelen.
\end{solution}

\begin{solution}{pb:b3:bacteriology:1}
\textbf{1.} $18$ generaties: $2^{18} = 2.6\times 10^{5}$ cellen,
$2.6\times 10^{3}$ per mL.
\textbf{2.} $2\times 10^{9}\times 100 = 2\times 10^{11}$ cellen $=
2^{37.5}$: na \qty{12.5}{h}; massa \qty{0.2}{g}.
\textbf{3.} $6\times 10^{27}/10^{-12} = 6\times 10^{39} = 2^{132}$
cellen: $132$ generaties, \qty{44}{h}. Voedingsstoffen, zuurstof en
afval leggen het binnen een dag stil; exponentiële groei duurt altijd
kort.
\textbf{4.} Cellen uit de \omterm{def:b3:bacteriology:phases}{stationaire fase} hebben hun ribosomen
afgebroken en hun stofwisselingsenzymen onderdrukt; zij moeten die
eerst weer opbouwen voordat zij zich delen. Exponentiële cellen komen
volledig uitgerust aan.
\textbf{5.} $\mu = \ln 2/\qty{0.333}{h} = \qty{2.08}{h^{-1}}$; bij $T =
\qty{60}{min}$ \qty{0.69}{h^{-1}}; na \qty{6}{h} $2^{6} = 64$ in plaats
van $2^{18}$: een factor $2^{12} \approx 4000$ minder.
\textbf{6.} $200/(0.1\times 10^{-6}) = 2\times 10^{9}$ levende cellen
per mL. De microscoop telt ook dode cellen en cellen die op de plaat
niet groeien, en een klontje geeft één kolonie.
\textbf{7.} $D = 0.5/1 = \qty{0.5}{h^{-1}}$; $T = \ln 2/0.5 =
\qty{1.39}{h} = \qty{83}{min}$.
\textbf{8.} $S^{*} = 0.01\times 0.5/0.5 = \qty{0.01}{g/L}$; $X^{*} =
0.5\times 1.99 \approx \qty{1}{g/L} = 10^{12}$ cellen per liter, $10^{9}$
per mL.
\textbf{9.} $0.5\times 10^{12} = 5\times 10^{11}$ cellen per uur; $720/1.39
\approx 520$ generaties per maand.
\textbf{10.} $D = 0.95$: $S^{*} = 0.01\times 0.95/0.05 = \qty{0.19}{g/L}$,
$X^{*} = 0.5\times 1.81 = \qty{0.9}{g/L}$. Bij $F = \qty{1.1}{L/h}$ is $D
> \mu_{\max}$: uitspoeling.
\textbf{11.} De mutant houdt $S^{*} = 0.01\times 0.5/0.7 = \qty{0.0071}{g/L}$;
bij die glucose groeit de oorspronkelijke stam met $1.0\times 0.0071/0.0171 =
\qty{0.42}{h^{-1}} < D$ en neemt zij af als $e^{-0.08t}$: in enkele
weken uitgespoeld. De mutant neemt het over.
\textbf{12.} $10^{-9}\times 10^{12} = 1000$ gunstige mutaties per
generatie: de aanpassing is doorlopend, met om de paar honderd
generaties een nieuwe voordelige kloon die het overneemt --- de
chemostaat is een evolutiemachine.
\textbf{13.} $kA_{\text{thr}}/p = 0.05\times 6\times 10^{12}/500 =
6\times 10^{8}$ per liter $= 6\times 10^{5}$ cellen per mL.
\textbf{14.} $10^{10}/6\times 10^{5} \approx 1.7\times 10^{4}$-voudig.
\textbf{15.} $20\times 6\times 10^{12}/500 = 2.4\times 10^{11}$ per liter,
$2.4\times 10^{8}$ per mL --- zes orden van grootte boven de $10^{2}$
per mL van de zee: donker.
\textbf{16.} Tienvoudig is $3.3$ verdubbelingen, \qty{6.6}{h}. Bij
$10^{9}$ per mL ligt de populatie nog $1700$ keer boven de drempel: het
licht blijft aan.
\textbf{17.} In zee ziet niemand het licht en voedt het geen gastheer:
een donkere mutant bespaart \qty{20}{\%} en groeit de lichtgevers
voorbij. In het orgaan selecteert de inktvis voor licht --- zij stoot
donkere profiteurs uit of voedt ze niet --- zodat de kosten een
woonplaats kopen.
\textbf{18.} Blokkeer de receptor of vernietig de \omterm{def:b3:bacteriology:quorum}{auto-inductor}
(quorumdemping): de bacteriën leven maar zetten hun toxinen nooit in,
en het afweersysteem ruimt hen op. Omdat gevoelige en ``resistente''
bacteriën even hard groeien --- er wordt niets gedood --- is het
selectieve voordeel van aan het middel ontsnappen klein, en zou de
resistentie zich trager moeten verspreiden dan tegen een dodend middel.
\textbf{19.} A: $10^{-8}\times 10^{9} = 10$; B: $100$; beide: $10^{-15}\times
10^{9} = 10^{-6}$.
\textbf{20.} $P_{0}(\text{A}) = e^{-10} = 5\times 10^{-5}$;
$P_{0}(\text{beide}) = e^{-10^{-6}} = 0.999999$.
\textbf{21.} $16 \to 8$: één halveringstijd, \qty{4}{h}; $16 \to 1$:
vier, \qty{16}{h}; in het venster van \qty{4}{h} tot \qty{16}{h}, dus
\qty{12}{h} per dosis.
\textbf{22.} Bij dosering elke \qty{12}{h}: in het venster van uur $4$
tot uur $12$, twee derde van de tijd; over tien dagen worden de
mutanten van de eerste stap tegen A --- tien bij aanvang ---
geselecteerd en faalt A alleen.
\textbf{23.} Voortdurend boven \qty{8}{mg/L}: doseer elke halveringstijd
(\qty{4}{h}), of geef elke \qty{12}{h} een dosis met een piek van
\qty{64}{mg/L} ($64 \to 8$ in drie halveringstijden), of een doorlopend
infuus. De prijs: de giftigheid van hoge pieken, of zes doses per dag
die patiënten niet volhouden.
\textbf{24.} $10^{-4}\times 10^{9} = 10^{5}$ persisters overleven wat de
middelen ook doen; zodra de behandeling stopt ontwaken zij en groeien
zij weer uit tot een volkomen gevoelige populatie --- een terugval. De
behandeling moet lang genoeg duren om hen bij het verlaten van de
sluimering te vangen, en daarom kost tuberculose zes maanden.
\textbf{25.} $10^{9}$ cellen per mL in de chemostaat; het licht gaat aan
bij $6\times 10^{5}$ cellen per mL; $P_{0} = 0.999999$ dat de kuur met
twee middelen geen resistente cel tegenkomt.
\end{solution}
